% =============================================================================
% Chapter 1 -- Introduction
% Handbook guidance:
%   * Frame the real-world problem in plain language.
%   * Provide enough background that a non-specialist can follow.
%   * State Aims & Requirements as a NUMBERED list. Subsequent chapters
%     refer back to (R1), (R2), ... when reporting how each was met.
% =============================================================================
\chapter{Introduction}
\label{ch:intro}

An autonomous agent that is deployed once and never updated is rare:
in any setting where the world changes, the agent must keep learning
after the laboratory training run is over. A household robot
encounters new room layouts; a recommender meets new content
categories; a game-playing agent is asked to handle a patch that
changed the level. The natural way to extend a trained agent is to
keep training it on whatever new data arrives. The catch is that a
neural network trained on the new data alone tends to overwrite its
old skills almost immediately --- a phenomenon known as
\emph{catastrophic forgetting} \citep{french1999catastrophic,
kirkpatrick2017ewc}. The field of \emph{continual learning} (CL)
studies practical remedies for this failure mode.

A second, more pragmatic constraint is rarely discussed in published
benchmarks but dominates real deployments: \emph{compute}. Continual
RL benchmarks today report headline numbers obtained on multi-GPU
clusters with hundreds of thousands of environment steps per task. A
student replicating these numbers on a laptop, or a practitioner
fine-tuning a deployed model in the field, will not have that budget.
This dissertation asks the natural follow-on question: \emph{do the
ranking and the qualitative behaviour of standard continual RL methods
hold when the per-task compute budget is reduced by an order of
magnitude?}

The vehicle for the study is the COOM benchmark of
\citet{tomilin2024coom}, a continual reinforcement learning suite of
ViZDoom \citep{kempka2016vizdoom} mini-games designed to expose
exactly the catastrophic forgetting problem this dissertation
investigates. The compute regime is fixed by what an Apple Silicon
laptop can deliver overnight; the methods compared are the three
regularisation-based baselines reported in the COOM paper --- EWC, L2
anchoring and MAS --- against the same unregularised SAC fine-tuning
control. The result is a negative one, and the dissertation argues
that the negative result is itself informative: the published ranking
\emph{inverts} when compute is constrained, and the inversion has a
specific, diagnosable cause.

\section{Background}
\label{sec:background}

\subsection{Catastrophic forgetting and the regularisation family}

Catastrophic forgetting was first reported in connectionist networks
in the late 1980s \citep{french1999catastrophic} and remains, three
decades later, the defining obstacle in continual learning. Among the
many remedies proposed, three families dominate the literature
\citep{parisi2019continual}: \emph{regularisation-based} methods,
which constrain how far parameters may drift from their previous-task
values; \emph{rehearsal-based} methods, which interleave a small
amount of stored or generated past data into current training
\citep{chaudhry2019agem, daniels2022clonex}; and
\emph{architecture-based} methods, which give each task a
dedicated subset of parameters \citep{mallya2018packnet}. This
dissertation focuses on the first family, because it is parameter
efficient (no replay buffer is stored), architecturally lightweight
(no parameter masking), and because the COOM paper reports it as the
strongest non-rehearsal family on the benchmark.

The canonical regularisation method is Elastic Weight Consolidation
(EWC) \citep{kirkpatrick2017ewc}. After training on task $t$, the
diagonal of the Fisher information matrix $F^{(t)}$ is computed at
the converged parameter vector $\theta^*_t$ and used to weight a
quadratic anchor on subsequent tasks:
\begin{equation}
    \mathcal{L}_{\text{EWC}}(\theta) =
        \mathcal{L}_t(\theta)
        + \lambda \sum_i F^{(t-1)}_i (\theta_i - \theta^{*}_{t-1, i})^2.
    \label{eq:ewc}
\end{equation}
The Fisher diagonal estimates how much the data log-likelihood
changes locally with each parameter; intuitively, parameters that
matter for the source task are pulled back hard, and parameters that
do not are left free for the new task. L2 anchoring is the same
penalty with $F$ replaced by the identity. Memory Aware Synapses
(MAS) \citep{aljundi2018mas} replaces $F$ with the squared $L_2$ norm
of the gradient of the model output, an unsupervised importance
estimate that does not need a label or a likelihood.

\subsection{Reinforcement learning, SAC, and the COOM benchmark}

In a reinforcement learning setting the supervised log-likelihood
underlying \cref{eq:ewc} is unavailable. The COOM implementation
adapts EWC by computing the Fisher diagonal as the squared Jacobian
of the actor logits and Q-value heads with respect to the trainable
parameters, averaged over a small batch of samples drawn from the
replay buffer at the end of each task. The base RL algorithm is Soft
Actor-Critic \citep{haarnoja2018sac}, an off-policy maximum-entropy
method whose stochastic actor and twin critic make it a natural fit
for continuous-state, discrete-action ViZDoom environments. The CL
training loop is inherited from the Continual World framework
\citep{wolczyk2021cwsac}: tasks are presented sequentially, with the
agent training only on the current task while regularisation
penalties anchor it to its parameters at the end of each previous
task.

The COOM benchmark itself \citep{tomilin2024coom} packages eight
distinct ViZDoom scenarios into Cross-Domain (CD) sequences, in which
the same scenario is presented with progressively different
visuals and dynamics, and Cross-Objective (CO) sequences, in which
each task is a different scenario with a different objective. The
benchmark reports Average Performance, Forgetting and Forward
Transfer for each method at $200{,}000$ training steps per task and
ten random seeds. The headline ranking (Section~6 of the COOM paper)
places PackNet first, followed by ClonEx-SAC, with the
regularisation family clustered behind: L2 at $0.64$, MAS at $0.56$,
EWC at $0.54$. Crucially, every regularisation method scores above
the unregularised Fine-Tuning control at $0.40$.

\subsection{Plasticity loss in deep RL}

Independent of the CL literature, the deep RL community has reported
that even a single-task RL agent can lose its ability to learn after
prolonged training, a phenomenon known as \emph{plasticity loss}
\citep{lyle2022plasticity, nikishin2022primacy}. The mechanism is
not yet fully understood, but the practical lesson is that early
optimisation noise in an RL run is qualitatively different from the
behaviour at convergence. This will become relevant in
\Cref{ch:eval}, where the failure mode of low-budget EWC is
diagnosed as a Fisher matrix that captures early-training
optimisation noise rather than a converged source-task optimum.

\section{Aims and Requirements}
\label{sec:aims}

This project has the following primary aim:
\begin{quote}
    To empirically characterise the behaviour of regularisation-based
    continual reinforcement learning methods (EWC, L2, MAS) when the
    per-task training budget is reduced by an order of magnitude relative
    to the published COOM benchmark.
\end{quote}

The aim decomposes into the following numbered requirements, used as
acceptance criteria throughout the rest of this dissertation.

\begin{enumerate}[label=\textbf{R\arabic*.},leftmargin=2.4em]
    \item \label{req:pipeline}
        Reproduce the COOM training pipeline locally on consumer
        hardware, end to end, for at least one task sequence and one
        regularisation method.
    \item \label{req:multimethod}
        Run a controlled comparison of Fine-Tuning, EWC, L2 and MAS on
        the same sequence with identical hyperparameters except the
        regularisation method.
    \item \label{req:variance}
        Report seed-level variance for at least one method, providing a
        bound on the inter-seed offset that could affect the comparison.
    \item \label{req:diagnose}
        Diagnose any observed failure modes, distinguishing between
        implementation bugs, hyperparameter mistuning, and
        budget-induced behaviour.
    \item \label{req:reproduce}
        Release a documented, single-command pipeline that reproduces
        every figure and table in the dissertation.
\end{enumerate}

\section{Dissertation Structure}
\label{sec:structure}

The remainder of this dissertation is organised as follows.
\Cref{ch:reqs} sets out the project requirements in stakeholder
terms. \Cref{ch:design} justifies the experimental design that meets
requirements~\ref{req:pipeline} and~\ref{req:multimethod}, including
the choice of task sequence, training budget and method set, and
records the alternatives that were considered and rejected.
\Cref{ch:impl} documents the technical realisation: the
Apple-Silicon CPU TensorFlow stack, the sequential experiment queue
that schedules five overnight runs, and the reproducibility layer
that satisfies requirement~\ref{req:reproduce}. \Cref{ch:eval}
reports results addressing requirements~\ref{req:variance}
and~\ref{req:diagnose}: average training success, per-task curves,
within-task drift, and the regularisation-loss diagnostic that
localises the failure mode. \Cref{ch:ethics} discusses project
ethics. \Cref{ch:concl} maps every finding back to the numbered
requirements and identifies four concrete directions for follow-up
work; \Cref{ch:bcs} closes with the BCS criteria checklist and a
critical first-person reflection.
